% Zufallsgrößen und Verteilungen – bank/zufallsgroessen-und-verteilungen/ (zusammenbau v0.4, Heft)
\einheitenkopf[t5]{Zufallsgrößen und Verteilungen}
\setcounter{aufgabe}{24}

% Anlauf (ohne Prüfkennung)
\begin{aufgabe}{Verteilung aufstellen – Anlauf}
\begin{teile}
\teil Zwei Münzen werden geworfen; $X$ zählt, wie oft Kopf fällt – welche Werte kann $X$ annehmen? \kreuz{$0, 1, 2$}\\ \kreuz{$1, 2$}\\ \kreuz{$0, 1, 2, 3$}
\teil Zwei Tetraederwürfel (Augenzahlen $1$ bis $4$) werden geworfen; $X$ ist die Augensumme. Fülle die Tabelle. \sachtabelle{lcccc}{$k$ & 2 & 3 & 4 & 5}{$P(X = k)$ & \leerzelle & \leerzelle & \leerzelle & \leerzelle}
\teil Ein Glücksrad mit den gleich großen Feldern $1$ und $3$ wird zweimal gedreht; ausgezahlt wird das Produkt in Euro, der Einsatz ist $2$ €. Zeige, dass der Gewinn (Auszahlung minus Einsatz) nur die Werte $-1$, $1$ und $7$ annimmt.
\end{teile}
\end{aufgabe}

\begin{aufgabe}{Verteilung aufstellen – Prüfungsaufgaben}
\begin{teile}
\teil Drei Briefe werden zufällig in drei adressierte Umschläge gesteckt; $X$ ist die Anzahl der richtig steckenden Briefe. Es gilt $P(X = 0) = \frac{1}{3}$ und $P(X = 1) = \frac{1}{2}$. Bestimme $P(X = 2)$ und $P(X = 3)$. \hfill (Abitur 2022 GK)
\teil Zwei schwarze und zwei weiße Socken werden zufällig auf zwei Schubladen verteilt, je zwei in jede; $X$ ist die Anzahl der Schubladen mit zwei verschiedenfarbigen Socken. Es gilt $P(X = 0) = \frac{1}{3}$. Bestimme $P(X = 1)$ und $P(X = 2)$. \hfill (Abitur 2022 GK)
\end{teile}
\end{aufgabe}
